\chapter{Hot Flushes (Hot Flashes)}

\section*{Definition}
Sudden episodes of intense heat sensation spreading over the upper body and face, often accompanied by sweating, palpitations, and anxiety, primarily associated with hormonal changes.

\section*{What Happens in Your Body}
Estrogen withdrawal (menopausal transition) alters brain's temperature control center thermoregulatory set-point, narrowing the thermoneutral zone so minor core temperature elevations trigger heat-loss responses (widening of blood vessels, sweating). Fluctuations in luteinizing hormone and norepinephrine also contribute to vasomotor instability.

\section*{Causes}
\begin{itemize}
  \item Menopause (natural or surgical)
  \item Perimenopause
  \item Medication side effects (SSRIs, tamoxifen, raloxifene)
  \item Thyroid disorders (hyperthyroidism)
  \item Cancer treatment (chemotherapy, radiation, hormone therapy)
  \item Orchidectomy (in men)
  \item Carcinoid syndrome
  \item Adrenal insufficiency
  \item Alcohol or caffeine triggers
  \item Anxiety or panic disorder
\end{itemize}

\section*{When to See a Doctor}
See your doctor if:
\begin{itemize}
  \item Symptoms are severe or getting worse over time
  \item Hot flushes with fever, night sweats with weight loss, or flushes interfering significantly with quality of life
  \item You have other concerning symptoms such as fever, weight loss, or bleeding
\end{itemize}

\section*{Related Symptoms}
\begin{itemize}
  \item Night sweats
  \item Palpitations
  \item Anxiety
  \item Sleep disturbance
  \item Mood swings
\end{itemize}
\textbf{Important:} If this symptom persists, your doctor will check for serious underlying causes.

\section*{Self-Care / Home Management}
\begin{itemize}
  \item Rest and avoid activities that make it worse.
  \item Maintain adequate hydration and a balanced diet.
  \item Over-the-counter remedies may provide symptomatic relief (consult a pharmacist or doctor).
  \item Monitor symptoms and seek professional advice if they persist or worsen.
  \item Avoid self-medicating with prescription medications without medical guidance.
\end{itemize}

\section*{How Your Doctor Will Evaluate You}
\begin{itemize}
  \item A thorough history and physical examination are the first steps in evaluation.
  \item Laboratory tests (blood work, urinalysis) may be ordered based on clinical suspicion.
  \item Imaging studies (X-ray, ultrasound, CT, MRI) may be indicated when structural pathology is suspected.
  \item Specialist referral is arranged when the diagnosis remains uncertain or requires specific expertise.
\end{itemize}

\section*{Treatment Options}
\begin{itemize}
  \item Treatment targets the underlying cause identified through diagnostic evaluation.
  \item Supportive care and symptom management are provided while awaiting definitive therapy.
  \item Medications, lifestyle modifications, and physical therapies may be prescribed as appropriate.
  \item Follow-up ensures treatment effectiveness and allows timely adjustments to the care plan.
\end{itemize}

\clearpage